\documentclass[12pt,a4paper]{article}
\usepackage[italian]{babel}
\usepackage[utf8]{inputenc}
\usepackage{amsmath, amssymb}
\usepackage{graphicx}
\usepackage{hyperref}
\usepackage{geometry}
\usepackage{listings}
\usepackage{longtable}
\geometry{margin=2.5cm}

\title{Clustering}
\author{Prof. Fedeli Massimo}
\date{}

\lstset{
	basicstyle=\ttfamily\small,
	frame=single,
	breaklines=true
}

\begin{document}
	
	\maketitle
	\tableofcontents
	\newpage
	
	\section{Introduzione al Clustering}
	
	Il clustering è una delle tecniche fondamentali del machine learning non supervisionato. A differenza dei modelli supervisionati, non si dispone di etichette, e l'obiettivo è individuare strutture latenti nei dati.
	
	\subsection{Perché è importante}
	
	\begin{itemize}
		\item Analisi esplorativa dei dati
		\item Segmentazione clienti
		\item Compressione dell'informazione
		\item Rilevazione anomalie
	\end{itemize}
	
	\subsection{Esempio concreto}
	
	Un'azienda di e-commerce vuole dividere i clienti in gruppi. I dati disponibili sono:
	\begin{itemize}
		\item età
		\item spesa media
		\item frequenza acquisti
	\end{itemize}
	
	Il clustering consente di individuare gruppi con comportamenti simili.
	
	\section{Preparazione dei dati}
	
	\subsection{Normalizzazione}
	
	\begin{lstlisting}[language=Python]
		from sklearn.preprocessing import StandardScaler
		X_scaled = StandardScaler().fit_transform(X)
	\end{lstlisting}
	
	\subsection{Perché è fondamentale}
	
	Senza normalizzazione, variabili con scale diverse dominano la distanza.
	
	\section{K-Means Approfondito}
	
	\subsection{Funzione obiettivo}
	
	\begin{equation}
		J = \sum_{i=1}^K \sum_{x \in C_i} ||x - \mu_i||^2
	\end{equation}
	
	\subsection{Algoritmo dettagliato}
	
	\begin{enumerate}
		\item scelta casuale centroidi
		\item assegnazione punti
		\item aggiornamento centroidi
		\item iterazione
	\end{enumerate}
	
	\subsection{Esempio numerico}
	
	Consideriamo 4 punti: (1,1), (1,2), (5,5), (6,5).
	
	Con K=2, i cluster convergono naturalmente in due gruppi distinti.
	
	\subsection{K-Means++}
	
	Tecnica migliorata per inizializzazione.
	
	\subsection{Caso in cui fallisce}
	
	\begin{itemize}
		\item cluster non sferici
		\item presenza di outlier
	\end{itemize}
	
	\subsection{Codice completo}
	
	\begin{lstlisting}[language=Python]
		from sklearn.cluster import KMeans
		
		kmeans = KMeans(n_clusters=3, init='k-means++')
		labels = kmeans.fit_predict(X_scaled)
	\end{lstlisting}
	
	\section{DBSCAN Approfondito}
	
	\subsection{Definizioni}
	
	\begin{itemize}
		\item punto core
		\item punto border
		\item rumore
	\end{itemize}
	
	\subsection{Intuizione}
	
	Cluster = regioni dense separate da zone sparse.
	
	\subsection{Esempio}
	
	DBSCAN riesce a separare cluster non lineari.
	
	\subsection{Scelta parametri}
	
	\begin{itemize}
		\item eps: raggio
		\item minPts: densità minima
	\end{itemize}
	
	\subsection{Codice}
	
	\begin{lstlisting}[language=Python]
		from sklearn.cluster import DBSCAN
		
		db = DBSCAN(eps=0.3, min_samples=5)
		labels = db.fit_predict(X_scaled)
	\end{lstlisting}
	
	\section{Clustering Gerarchico}
	
	\subsection{Agglomerativo}
	
	Parte da singoli punti.
	
	\subsection{Divisivo}
	
	Parte da un unico cluster.
	
	\subsection{Linkage}
	
	\begin{itemize}
		\item single
		\item complete
		\item average
	\end{itemize}
	
	\subsection{Codice}
	
	\begin{lstlisting}[language=Python]
		from sklearn.cluster import AgglomerativeClustering
		
		model = AgglomerativeClustering(n_clusters=3, linkage='ward')
		labels = model.fit_predict(X_scaled)
	\end{lstlisting}
	
	\section{Gaussian Mixture Models}
	
	\subsection{Modello}
	
	\begin{equation}
		p(x) = \sum_{k=1}^K \pi_k \mathcal{N}(x|\mu_k,\Sigma_k)
	\end{equation}
	
	\subsection{Algoritmo EM}
	
	\begin{itemize}
		\item E-step
		\item M-step
	\end{itemize}
	
	\subsection{Codice}
	
	\begin{lstlisting}[language=Python]
		from sklearn.mixture import GaussianMixture
		
		gmm = GaussianMixture(n_components=3)
		labels = gmm.fit_predict(X_scaled)
	\end{lstlisting}
	
	\section{Metriche di valutazione}
	
	\subsection{Silhouette}
	
	\begin{lstlisting}[language=Python]
		from sklearn.metrics import silhouette_score
		
		score = silhouette_score(X_scaled, labels)
	\end{lstlisting}
	
	\subsection{Davies-Bouldin}
	
	Valuta separazione tra cluster.
	
	\section{Confronto algoritmi}
	
	\begin{longtable}{|c|c|c|c|}
		\hline
		Algoritmo & Tipo & Parametri & Outlier \\
		\hline
		K-Means & Partizionale & K & No \\
		DBSCAN & Densità & eps & Sì \\
		Gerarchico & Gerarchico & linkage & Limitato \\
		GMM & Probabilistico & K & No \\
		\hline
	\end{longtable}
	
	\section{Caso studio completo}
	
	\subsection{Scenario}
	
	Segmentazione clienti retail.
	
	\subsection{Pipeline}
	
	\begin{enumerate}
		\item pulizia dati
		\item scaling
		\item clustering
		\item analisi risultati
	\end{enumerate}
	
	\subsection{Codice completo}
	
	\begin{lstlisting}[language=Python]
		from sklearn.preprocessing import StandardScaler
		from sklearn.cluster import KMeans
		
		X_scaled = StandardScaler().fit_transform(X)
		
		kmeans = KMeans(n_clusters=4)
		labels = kmeans.fit_predict(X_scaled)
	\end{lstlisting}
	
	\section{Errori comuni}
	
	\begin{itemize}
		\item non normalizzare
		\item scegliere K arbitrariamente
		\item ignorare outlier
	\end{itemize}
	
	\section{Conclusioni}
	
	Il clustering è uno strumento potente ma richiede attenzione nella preparazione dei dati e nella scelta dell'algoritmo.
	
\end{document}